\section{The \texttt{Shadow\_output} McXtrace Component}
Release: McXtrace 0.1

Write x-ray state parameters to SHADOW x-ray event file.

\subsection*{Identification}
\begin{itemize}
  \item \textbf{Author:} Andrea Prodi
  \item \textbf{Origin:} Risoe/ILL
  \item \textbf{Date:} November 21, 2011
\end{itemize}

\subsection*{Description}
Detector-like component writing x-ray state parameters to a SHADOW x-ray file. Used to interface McXtrace components or simulations into SHADOW. Each photon is 104 bytes.

Note that when standard output is used, as is the default, no monitors or other components that produce terminal output must be used, or the x-ray output from this component will become corrupted.

Example: Shadow\_output(file="MySource.vit", bufsize = 10000, progress = 1)

\subsection*{Input parameters}
Parameters in \textbf{boldface} are required; the others are optional.

\begin{longtable}{p{0.22\textwidth}p{0.12\textwidth}p{0.46\textwidth}p{0.14\textwidth}}
\toprule
\textbf{Name} & \textbf{Unit} & \textbf{Description} & \textbf{Default} \\
\midrule
\endhead
file & string & Filename of x-ray file to write. Default is standard output & "" \\
bufsize & records & Size of x-ray output buffer & 1000 \\
progress & flag & If not zero, output dots as progress indicator & 0 \\
\bottomrule
\end{longtable}

\subsection*{Links}
\begin{itemize}
  \item Component source code found in file \texttt{Shadow\_output.comp}.
\end{itemize}
\IfFileExists{misc/Shadow_output_static.tex}{\input{misc/Shadow_output_static.tex}}{}